\section{The acquired action-occupation theorem}\label{sec:central}
For $C=(r_1,\ldots,r_m)\in H^m$ put $D_C(x)=\min_j\norm{x-r_j}^2$. Define
\begin{align}
 \Psi_m(N;c)=\min\ &\sum_{a\in\calA}\int D_C(x)h_a(x)\nu_a(dx),\label{eq:profile}\\
 &C\in H^m,\quad x_a\ge0,\quad 0\le h_a\le x_a\quad\nu_a\text{-a.e.},\notag\\
 &\sum_a\int h_a\,d\nu_a=1,\qquad \sum_a c_ax_a\le N.\notag
\end{align}
An infeasible minimum is $+\infty$. Here $x_a$ is an expected acquisition count, and $h_a\nu_a$ is an \emph{unnormalized} terminal submeasure. Repeated letters $x$ in the integration variable and $x_a$ in the counts have different types. We write $\Psi_m(N)$ when all costs are one.

\begin{theorem}[Causal acquired-geometry and acquisition transfer]\label{thm:central}
Let the experiment and task satisfy Section~\ref{sec:experiment}, and let $m\ge1$, $M\ge2$.

\textup{(i)} Under task regeneration, for every real $N\ge\min_a c_a$ the minima exist and
\begin{equation}\label{eq:exact}
 \calR_\cp(N,m;c)-B=\Psi_m(N;c),\qquad
 \calR_\aut(N,M;c)-B=\Psi_{M-1}(N;c).
\end{equation}
Every feasible $(x_a,h_a)_a$ in \eqref{eq:profile} is realized, with its exact terminal submeasures and costs, by a homogeneous observer with one live label and at most $m$ terminal labels. An optimal observer can be chosen to use at most two acquisition actions.

\textup{(ii)} For fixed centers the profile has the dual representation
\begin{equation}\label{eq:dual}
 V_C(N;c)=\sup_{\lambda\ge0,\ z\in\R}
 \left\{z-\lambda N:
       \int(z-D_C)_+\,d\nu_a\le\lambda c_a\quad\text{for every }a\right\},
 \qquad \Psi_m=\min_{C\in H^m}V_C.
\end{equation}
For $N>\min_a c_a$ the dual is attained. At the endpoint its supremum remains exact, without an assertion of attainment. The same pair $(z,\lambda)$ is used for every action. No exchange of the center minimum and the dual supremum is asserted.

\textup{(iii)} Under the conditional envelopes \eqref{eq:envelopes}, let $\kappa=\max_a\overline b_a/\underline b_a$. Then the common-program optima obey
\begin{align}
 \Psi_m(N;c/\overline b)&\le\calR_\cp(N,m;c)-B,
 &\calR_\cp(N,m;c)-B&\le\kappa\Psi_m(N;c/\underline b),\label{eq:cp-envelope}\\
 \Psi_{M-1}(N;c/\overline b)&\le\calR_\aut(N,M;c)-B,
 &\calR_\aut(N,M;c)-B&\le\kappa\Psi_{M-1}(N;c/\underline b).
 \label{eq:aut-envelope}
\end{align}
The right-hand bounds assert an implementation only when the respective profile is feasible. Its program is parameter independent and still has one live label.

\textup{(iv)} Consider two known regenerative response families on the same compact task space, with unit acquisition costs and the same protected audit mean and constant variance. If $\max_a\norm{\nu_a-\nu'_a}_\TV\le\eta$, and $0\le D_C\le L$ for all projected decoders, then
\begin{equation}\label{eq:linear}
 |\Psi_m(N)-\Psi'_m(N)|\le\min\{L,2LN\eta\}\qquad(N\ge1).
\end{equation}
The comparison transfers an optimal one-live-state program to a program with the \emph{same} state and mean budgets. It is not an inference that total variation controls arbitrary stopping means.
\end{theorem}

The theorem optimizes over all nonanticipating acquisition strategies, not merely over a common fixed controller. Its exact part is also a sufficiency statement: the scalar mass inequalities characterize the entire attainable selected law on this class. Task regeneration is essential for this characterization. Under \eqref{eq:envelopes} the theorem asserts the displayed matching comparisons, not equality of selected laws. The Gaussian result of Section~\ref{sec:gaussian} shows how an informative experiment outside this exact class is instead handled by the raw occupation equations.

\begin{corollary}[Regular acquired supports]\label{cor:regular}
Suppose each $\nu_a$ has support $K_a$ of diameter at most a fixed constant and, for fixed $c_0,C_0,r_0>0$,
\begin{equation}\label{eq:regular}
 c_0r^{s_a}\le\nu_a(B(x,r))\le C_0r^{s_a}
 \quad(x\in K_a,\ 0<r\le r_0),\qquad
 0<s_-\le s_a\le s_+<\infty.
\end{equation}
For unit costs and uniformly bounded envelope constants, for every $N\ge1$,
\begin{align}
 \calR_\cp(N,m)-B&\asymp\min_a(Nm)^{-2/s_a},\label{eq:regular-cp}\\
 \calR_\aut(N,M)-B&\asymp\min_a\{N(M-1)\}^{-2/s_a}.
 \label{eq:regular-auto}
\end{align}
Constants depend only on the displayed regularity, diameter and envelope data and the finite action family, not on $N,m,M$. These exponents belong to actual one-call supports. They are not the dimension of a chosen physical or parameter carrier.
\end{corollary}

For nonregular or mixed laws it is \eqref{eq:profile}, not a single exponent, that remains the conclusion. The exact profile allows actions to have different atomic masses and different scales; mixtures may be useful at finite budgets even when one action eventually has the smallest regular exponent. No optimal-control theorem is inferred for actions outside the declared legal family.
